

\chapter{Appendix A}

\renewcommand{\thefigure}{A.\arabic{figure}}
\section{Risk Analysis}
\label{sec:risks}

This lists the main risks, with their severity and likelihood, whilst also considering any mitigation. A small table was previously mentioned in Chapter~\ref{sec:feas-risks}, but this is the complete version of the risk analysis.

\begin{table}[!ht]
  \centering
  \footnotesize
  \renewcommand{\arraystretch}{1.2}
  \caption{Risk analysis.}
  \label{tab:risks}
  \begin{tabular}{|p{0.27\linewidth}|p{0.09\linewidth}|p{0.09\linewidth}|p{0.45\linewidth}|}
    \hline
    \textbf{Risk} & \textbf{Severity} & \textbf{Likelihood} & \textbf{Mitigation} \\
    \hline
    A model on rendered images might not work on real images. & \High & \Med & Use the ECG-Image-Kit to generate a large number of images. Added augmentation of real artefacts to these images to make them more realistic. Report the gap as a result. \\
    \hline
    Trace detection falls apart on very poor quality images. & \Med & \Med & Detection is evaluated separately on real images. If it is too unreliable for analysis, the CNN uses the image alone. The failure is reported. \\
    \hline
    Some classes have too few examples to train a model. & \Med & \High &
    Results are reported per class. Some classes like AFLT and AFIB are combined into a single class. \\
    \hline
    Class imbalances cause the model to be biased towards the majority class. & \Med & \High &
    Use class weights, to balance the loss function. Use metrics that are aware of imbalances, (like sensitivity, specificity, AUROC). Thresholds are adjusted to account for the imbalance. \\
    \hline
    ECG-Image-Database has very few images for some classes. & \Med & \High &
    Only classes with enough images are evaluated. The rest are reported as a limitation. \\
    \hline
    Labels are subjective and may be inconsistent. & \Med & \Med &
    Validation and test folds that were checked by multiple annotators are used. The model is then evaluated on these folds. \\
    \hline
    Heatmaps can be misleading. & \Med & \Med & Heatmaps will be checked against known cases and shortcuts. The grid is masked if its needed. The heatmaps are evaluated in the user study. \\
    \hline
    The CNN underperforms. & \Med & \Med & A pretrained model might be used for comparison. \\
    \hline
    Training is too slow. & \Low & \Low & Use smaller models. This is to reduce the cost. \\
    \hline
  \end{tabular}
\end{table}

\begin{table}[!ht]
  \centering
  \footnotesize
  \renewcommand{\arraystretch}{1.2}
  \caption{Project and Ethical Risks.}
  \label{tab:project-risks}
  \begin{tabular}{|p{0.27\linewidth}|p{0.09\linewidth}|p{0.09\linewidth}|p{0.45\linewidth}|}
    \hline
    \textbf{Risk} & \textbf{Severity} & \textbf{Likelihood} & \textbf{Mitigation} \\
    \hline
    The scope is too large for the time available. & \Med & \Med & Requirements are prioritised with MoSCoW. This means that most important requirements are implemented first. In addition, the interface can be reduced to a single view if time is short. \\
    \hline
    There are not enough participants for the user study. & \Med & \Med & Recruitment of users starts early. If this fails, then an expert review replaces the study. \\
    \hline
    Clinicians are not able to take part. & \Low & \High & Requirements will be based on published studies from clinicians. Clinician testing is described at future work, and in addition the limitation is stated. \\
    \hline
    Software problems. & \Low & \Med & Use version control. Also, a virtual environment can be used with specific versions. This will have a recorded list of requirements. \\
    \hline
    Loss of code or data & \High & \Low & Keep remote copies. Also keep regular backups of results. \\
    \hline
    Work goes beyond the scope & \High & \Low & The work will stay in the approved scope. User data is destroyed upon completion of the project. Images are processed locally. \\
    \hline
    A dataset or tool update may change its availability or license & \Low & \Low & Local copies of the datasets are kept. Licenses are checked and cited. \\
    \hline
    Author might not be able to complete the project due to illness, or other personal reasons. & \High & \Low & The project is planned with deadlines, and the supervisor is kept informed of the progress. \\
    \hline
  \end{tabular}
\end{table}

\section{PLES}
\label{app:ples}


This section considers the professional, legal, ethical and social issues that this project raises. The project is a research prototype, that analyses ECG images and presents results to students and other users. It is not supposed to be for clinical use.
\subsection{Professional Issues}
The project follows the BCS Code of Conduct. These are the principles that apply most directly:

\begin{itemize}
  \item Public interest means that results are reported with honesty, even if they do not support my claim fully.  
  \item All datasets, tools, and papers used are cited.
  \item The preliminary results are described as pilot results.
  \item The use of generative AI is declared in the GenAI summary.
  \item All code is kept under version control
  \item Coding environment and its data is documented, so results can be reproduced.
\end{itemize}

\subsection{Legal Issues}
As ECG data is considered medical data, the UK Data Protection Law treats this as special category data. The training data comes from public, anonymised datasets, so the project does not process patient records that can be identified. The user study will still collect questionnaire responses, but they will be kept anonymous, and deleted upon project completion. Images uploaded to the interface are processed on the local machine, and will not be stored beyond the session.

Software that is intended to support clinically can be regulated as a medical device. This project is an academic prototype, and so will not be placed on markets or used on patients. This has been stated as such. Any later clinical use would need regulation, a clinical safety assessment, and validation on more data than the project uses.
PTB-XL and ECG-Image-Database are released under CC BY 4.0, whilst ECG-Image-Kit is under BSD 3-Clause License. They are cited whenever they are used, and the images generated by the toolkit are used under these terms. Figures from other sources are reproduced with a citation, and figures that the author produces are labelled as such. Copyright in the code, and the report, belong to the author under the university's regulations. 

\subsection{Ethical Issues}
The project will have ethical approval. Participants in the user study will give informed consent, in addition to be able to withdraw at any time, and are also not asked personal or health information.

The main risk of a tool that flags abnormalities is that users may rely on it. False negatives or positives can cause unecessary worry. Automation bias is documented in computer-interpreted ECGs. Inexperienced
users may accept an automated reading without question,
and explanations can raise the trust of novices more than experts. The project therefore reports a calibrated confidence for each
class so that any uncertainty is shown. It also checks heatmaps against known cases, and presents the
tool as a learning and research aid whose output has to be checked by a person.
Users are told how the model was trained, and where its weaknesses are.

\subsection{Social Issues}
As PTB-XL was collected 40 years ago, patient recording are limited. A model trained on this may perform differently for other populations, and so the project is not able to test this. The report does not claim that the results would hold on another system.

The tool could help students learn about ECG interpretation, and systems that read ECG images could widen the access to automated analysis, where only paper records exist. In addition, the interface should be usable by people with different needs, and any explanations should be clear and simple. However, users that over-rely on the tool may learn less information. Therefore, it is designed to explain the reasoning, to support learning.

For the environment, training deep models use a lot of energy. The project limits this using a small model, and pretrained weight if possible.
\section{Gantt Chart}
\label{sec:gantt}
\begin{figure}[!ht]
  \centering
  \resizebox{\linewidth}{!}{%
  \begin{ganttchart}[
      hgrid,                          % horizontal grid lines
      vgrid,                          % vertical grid lines
      x unit=0.55cm,                  % width of one week
      y unit chart=0.6cm,             % height of one row
      time slot format=simple,        % plain numbers instead of dates
      title label font=\scriptsize,
      bar/.append style={fill=blue!30},
      group/.append style={fill=black!70},
      milestone/.append style={fill=red!60},
      link/.style={-latex, thick, black!60}
    ]{1}{24}                          % weeks 1 to 24

    \gantttitle{Week}{24} \\
    \gantttitlelist{1,...,24}{1} \\

    \ganttgroup{Research}{1}{6} \\
    \ganttbar[name=lit]{Literature review}{1}{4} \\
    \ganttbar[name=req]{Requirements}{3}{6} \\

    \ganttgroup{Implementation}{6}{18} \\
    \ganttbar[name=data]{Dataset generation}{6}{9} \\
    \ganttbar[name=cnn]{CNN training}{9}{15} \\
    \ganttbar[name=gui]{GUI development}{12}{18} \\
    \ganttmilestone[name=proto]{Prototype complete}{15} \\

    \ganttgroup{Evaluation}{16}{22} \\
    \ganttbar[name=study]{User study}{17}{20} \\
    \ganttbar[name=write]{Dissertation writing}{16}{24} \\
    \ganttmilestone{Submission}{24}

    \ganttlink{lit}{req}
    \ganttlink{data}{cnn}
    \ganttlink{cnn}{proto}
    \ganttlink{gui}{study}
  \end{ganttchart}%
  }
  \caption{Project plan.}
  \label{fig:gantt}
\end{figure}
\section{GenAI Summary}
Generative AI was used for LaTex debugging, and grammatical as well as spelling errors. It was also used to help with formatting certain parts of the structure, but all the written body is in my own words. Every reference and citation I have found without the use of generative ai.

\section{Glossary}
\section*{Scripts}

Supplementary material...



\section*{GitHub Repository}

Project planning is core aspect of any dissertation (especially for the initial groundwork).  Key tool is Git/Github

